% ---------------------------------------------------------------------------
% Chapter: Installation
% ---------------------------------------------------------------------------
\chapter{Installation}\label{chap:install}
% Long monospaced identifiers leave little room for line breaking; allow some
% extra inter-word stretch in this chapter (the group keeps it local).
\begingroup\setlength{\emergencystretch}{3em}% chapter-local

My Tool Box is a standard Python project described by a single
\code{pyproject.toml} file. It is installed from its source repository with
\code{pip}; heavy dependencies are optional and grouped into \emph{extras}.

\section{Requirements}\label{sec:install-requirements}

\begin{itemize}
  \item Python 3.9 or newer. The package is tested with CPython 3.9 to 3.13;
    continuous integration runs the test suite on 3.9, 3.10, 3.11 and 3.12.
  \item Any operating system supported by the dependencies (Linux, macOS and
    Windows). Development and continuous integration use Linux.
  \item The core dependencies below, installed automatically by \code{pip}.
\end{itemize}

\begin{table}[htbp]
  \centering
  \small
  \caption{Core dependencies (always installed).}
  \label{tab:install-core}
  \begin{tabular}{@{}lll@{}}
    \toprule
    Package & Minimum version & Used for \\
    \midrule
    \pkg{numpy}      & 1.23 & all numerical code \\
    \pkg{scipy}      & 1.9  & AJLATT estimation, plotkit smoothing \\
    \pkg{matplotlib} & 3.6  & environment renderers, plotkit \\
    \pkg{gymnasium}  & 1.0  & environment API, registration, vector environments \\
    \pkg{pyyaml}     & 6.0  & AJLATT map files and map specifications \\
    \bottomrule
  \end{tabular}
\end{table}

\section{Installing from the repository}\label{sec:install-pip}

The recommended procedure creates a virtual environment, clones the
repository and installs the package with all optional features:

\begin{shellcode}
git clone https://github.com/EricDMW/My_Tool_Box.git
cd My_Tool_Box
python -m venv .venv
source .venv/bin/activate          # Windows: .venv\Scripts\activate
python -m pip install --upgrade pip
pip install ".[all]"
\end{shellcode}

\code{pip} can also install directly from the Git URL without a local clone,
which is convenient for a project that only uses the package:

\begin{shellcode}
pip install "my-tool-box[all] @ git+https://github.com/EricDMW/My_Tool_Box.git"
\end{shellcode}

With conda, create the environment first and then use \code{pip} inside it:

\begin{shellcode}
conda create -n mytoolbox python=3.11
conda activate mytoolbox
conda install tk                   # only for Tk windows (see below)
pip install ".[all]"
\end{shellcode}

\subsection{Optional extras}\label{sec:install-extras}

The extras of Table~\ref{tab:install-extras} are selected in square brackets
and can be combined, for example \code{pip install ".[pistonball,video]"}. The
quotes prevent shells such as \code{zsh} from interpreting the brackets.

\begin{table}[htbp]
  \centering
  \small
  \caption{Optional extras defined in \code{pyproject.toml}.}
  \label{tab:install-extras}
  \begin{tabularx}{\textwidth}{@{}l>{\raggedright\arraybackslash}p{4.3cm}>{\raggedright\arraybackslash}X@{}}
    \toprule
    Extra & Installs & Enables \\
    \midrule
    \code{pistonball} & \pkg{pygame} $\geq$ 2.1, \pkg{pymunk} $\geq$ 6.4 &
      \code{PistonballEnv} (\pkg{pymunk} for the physics, \pkg{pygame} for
      rendering and keyboard control) \\
    \code{torch} & \pkg{torch} $\geq$ 2.0 &
      \code{KuramotoOscillatorEnvTorch}, \pkg{toolkit.neural\_toolkit} \\
    \code{video} & \pkg{imageio} $\geq$ 2.20, \pkg{imageio-ffmpeg} $\geq$ 0.4 &
      MP4, WebM, AVI, MOV and MKV export in \code{save\_animation} and
      \code{record\_episode} (GIF export needs no extra) \\
    \code{all} & the three extras above & every feature of the package \\
    \code{dev} & \pkg{pytest} $\geq$ 7, \pkg{pytest-cov} $\geq$ 4, \pkg{ruff} $\geq$ 0.5 &
      running the test suite and the linter; combine with \code{all} \\
    \bottomrule
  \end{tabularx}
\end{table}

\subsection{PyTorch}\label{sec:install-torch}

The \code{torch} extra installs the default PyTorch wheel for the platform.
To choose a specific build, install PyTorch first, following the selector at
\url{https://pytorch.org/get-started/}, and then install the package; \code{pip}
keeps the PyTorch that is already present. For example, a CPU-only build
(much smaller, and the one used in continuous integration) is obtained with

\begin{shellcode}
pip install torch --index-url https://download.pytorch.org/whl/cpu
pip install ".[all]"
\end{shellcode}

GPU use is covered in Section~\ref{sec:trouble-torch}.

\subsection{PettingZoo}\label{sec:install-pettingzoo}

The adapter \code{env\_lib.wrappers.to\_parallel} exposes every environment
through the PettingZoo Parallel API (Section~\ref{sec:workflow-wrappers}).
PettingZoo is neither a dependency nor an extra: the adapter implements the
protocol itself. When PettingZoo is installed, the adapter subclasses
\code{pettingzoo.utils.env.ParallelEnv}, so \code{isinstance} checks and
PettingZoo's own utilities, such as \code{parallel\_api\_test}, accept it.
Install it only if your multi-agent library needs it:

\begin{shellcode}
pip install pettingzoo
python -c "import env_lib.wrappers as w; print(w.PETTINGZOO_AVAILABLE)"
\end{shellcode}

\subsection{Tk}\label{sec:install-tk}

Tk is not installable with \code{pip}. It is needed only for the
\pkg{parakit} editor window (Section~\ref{sec:parakit-gui}) and for
\code{render\_mode="human"} with matplotlib's TkAgg backend; everything else,
including off-screen rendering, works without it. Install it with the system
package manager (\code{sudo apt-get install python3-tk} on Debian and Ubuntu,
\code{sudo dnf install python3-tkinter} on Fedora, \code{brew install
python-tk} with Homebrew on macOS) or with \code{conda install tk}. The
official Python installers for Windows and macOS already include Tk.

\section{Development installation}\label{sec:install-dev}

Contributors install the package in \emph{editable} mode, so that changes to
the sources under \code{src/} take effect without reinstalling, together with
the development tools:

\begin{shellcode}
git clone https://github.com/EricDMW/My_Tool_Box.git
cd My_Tool_Box
python -m venv .venv && source .venv/bin/activate
pip install -e ".[all,dev]"
\end{shellcode}

The console scripts and the version number are generated at installation
time; after changing \code{pyproject.toml} (for example adding an entry
point), run the \code{pip install -e} command again. The testing and linting
workflow is described in Section~\ref{sec:examples-contributing}.

\section{Verifying the installation}\label{sec:install-verify}

The following commands check, from quickest to most thorough, that the
package is installed correctly:

\begin{shellcode}
# 1. Versions of both import packages
python -c "import env_lib, toolkit; print(env_lib.__version__, toolkit.__version__)"

# 2. Table of the registered environments (the env-lib console script)
env-lib list
python -m env_lib list            # the same without the console script

# 3. Create every environment and take a few random steps
python examples/quickstart.py --steps 5

# 4. Full test suite (requires the dev extra)
pytest -q
\end{shellcode}

The first command prints \code{1.1.0 1.1.0}. \code{env-lib list} creates every
environment once and prints one row per registered id (18 rows) with its
family, number of agents, observation shape, action type and shape, whether
\code{env\_lib.make\_vec} uses a native batched implementation, and the
optional dependency it requires; an environment whose extra is not installed
is marked \code{(missing)} instead of failing. \code{env-lib list -{}-no-inspect}
skips creating the environments. The command line is described in
Section~\ref{sec:workflow-cli}. The quick-start script prints one line per
environment and skips, with a message, environments whose optional
dependencies are missing. The test suite skips tests of missing optional
dependencies in the same way.

The installation provides four console scripts:

\begin{center}
  \small
  \begin{tabular}{@{}ll@{}}
    \toprule
    Command & Purpose \\
    \midrule
    \code{env-lib}          & list, run, evaluate and benchmark the environments (Section~\ref{sec:workflow-cli}) \\
    \code{ajlatt-build-map} & build an AJLATT map from a YAML specification (Chapter~\ref{chap:ajlatt}) \\
    \code{ajlatt-plot-map}  & list and plot the bundled AJLATT maps (Chapter~\ref{chap:ajlatt}) \\
    \code{plotkit-gallery}  & gallery of plotkit figures (Section~\ref{sec:plotkit-cli}) \\
    \bottomrule
  \end{tabular}
\end{center}

\section{Optional dependency matrix}\label{sec:install-matrix}

Table~\ref{tab:install-matrix} summarises which features need which optional
packages. Because the package imports optional dependencies lazily, a missing
dependency only affects the feature that needs it; the error message names the
extra to install.

\begin{table}[htbp]
  \centering
  \small
  \caption{Features and their optional requirements.}
  \label{tab:install-matrix}
  \begin{tabularx}{\textwidth}{@{}>{\raggedright\arraybackslash}X>{\raggedright\arraybackslash}p{3.8cm}>{\raggedright\arraybackslash}p{2.5cm}@{}}
    \toprule
    Feature & Requires & Provided by \\
    \midrule
    NumPy environments (Kuramoto, LineMsg, WirelessComm, Consensus, PowerGrid,
      Platoon, AJLATT),
      \code{rgb\_array} rendering, GIF export & core only & core \\
    \code{PistonballEnv} & \pkg{pymunk} & \code{pistonball} \\
    Pistonball rendering and \code{ManualPolicy} & \pkg{pygame} & \code{pistonball} \\
    \code{KuramotoOscillatorEnvTorch} & \pkg{torch} & \code{torch} \\
    \pkg{toolkit.neural\_toolkit} & \pkg{torch} & \code{torch} \\
    MP4 and other video formats & \pkg{imageio}, \pkg{imageio-ffmpeg} & \code{video} \\
    \pkg{toolkit.plotkit} & core only & core \\
    Plotting pandas, PyTorch, TensorFlow or JAX data & the respective library & not required \\
    \pkg{toolkit.parakit} headless API & standard library & core \\
    \code{env-lib} command line, catalogue, vector environments, baselines,
      evaluation & core only & core \\
    \code{to\_parallel} as a \code{pettingzoo.utils.env.ParallelEnv} subclass &
      \pkg{pettingzoo} (the adapter works without it) & not required \\
    \pkg{parakit} editor window & Tk and a display & system package \\
    \code{render\_mode="human"} (matplotlib environments) &
      an interactive backend (TkAgg, QtAgg, \ldots) and a display & system package or \pkg{PyQt}/\pkg{PySide} \\
    Test suite and linting & \pkg{pytest}, \pkg{pytest-cov}, \pkg{ruff} & \code{dev} \\
    \bottomrule
  \end{tabularx}
\end{table}

\section{Upgrading and uninstalling}\label{sec:install-upgrade}

To upgrade a clone, pull the latest sources and reinstall:
\code{git pull} followed by \code{pip install ".[all]"} (or
\code{pip install -e ".[all,dev]"} for an editable installation).
Reinstalling also runs again the step that generates the console scripts, so
an editable installation made with version 1.0 gains the \code{env-lib}
command only after this step. Version 1.1 requires Gymnasium 1.0 or newer;
\code{pip} upgrades an older Gymnasium automatically, which may affect other
projects in the same environment (Section~\ref{sec:migration-11}).
\code{pip uninstall my-tool-box} removes the package. Users of versions before
1.0, which were installed from two separate project folders, should first
uninstall the old distributions and then follow Appendix~\ref{app:migration}.

\section{Building this manual}\label{sec:install-manual}

The manual is written in LaTeX and lives in \code{docs/manual/}. Building it
requires a TeX distribution that provides \code{pdflatex} and \code{latexmk},
such as TeX Live or MacTeX. On Debian and Ubuntu the packages
\code{latexmk}, \code{texlive-latex-extra}, \code{texlive-fonts-recommended}
and \code{lmodern} are sufficient. The
document uses only standard packages (\pkg{listings}, \pkg{booktabs},
\pkg{tabularx}, \pkg{hyperref}, \pkg{fancyhdr} and similar).

\begin{shellcode}
cd docs/manual
./build.sh                                   # writes output/main.pdf
# equivalently:
latexmk -pdf -interaction=nonstopmode -halt-on-error -outdir=output main.tex
latexmk -C -outdir=output                    # remove build products
\end{shellcode}

\code{latexmk} runs \code{pdflatex} as often as needed to resolve cross
references. All build products are written to \code{docs/manual/output/},
which is ignored by Git. The file \code{docs/manual/README.md} describes the
source layout of the manual.

\par\endgroup
